\chapter{Markets and Microstructure}\label{ch:iv2:markets-and-microstructure}

\begin{ivtrains}
Reading an order book, what the spread pays for, and what your own trading does to the price: queues, adverse selection, impact, TWAP, VWAP and optimal liquidation, short selling, and the rules that shape all of it. Each stem gives its numbers. The rules quoted are US ones unless a stem says otherwise.
\end{ivtrains}

\section{Reading the book}

\begin{question}[$\star$ \iqroles{QT, Trader, QR, Dev} \iqfirm{market maker, hedge fund} \iqlevel{junior}]\label{iq:iv2:markets-and-microstructure:1}
An order book shows bids of 300 shares at 99.98 and 500 at 99.97, and offers of 100 shares at 100.02, 400 at 100.03 and 1\,000 at 100.05. You send a market order to buy 300 shares. What is your average price, and what do you pay above the mid-price, per share? What does the book's best offer and mid-price become?
\end{question}

\begin{question}[$\star$ \iqroles{QT, Trader, QR} \iqfirm{market maker, hedge fund} \iqlevel{junior}]\label{iq:iv2:markets-and-microstructure:2}
The best bid is 100.00 for 900 shares, and the best offer is 100.02 for 100 shares. What is the mid-price, and what is the size-weighted microprice, which weights each side's price by the size on the other side? Which way would you expect the next move to be, and why?
\end{question}

\begin{question}[$\star$ \iqroles{QT, Dev, QD, Trader} \iqfirm{market maker, hedge fund} \iqlevel{junior}]\label{iq:iv2:markets-and-microstructure:3}
At one price, three resting sell orders arrive in this order: $A$ for 100 shares, $B$ for 200 and $C$ for 50. A buy order for 250 shares arrives. How is it shared under price-time priority? How is it shared under pure pro-rata, where each order fills in proportion to its size? Which rule rewards speed, and which rewards size?
\end{question}

\begin{question}[$\star$ \iqroles{QT, Trader, QR} \iqfirm{market maker, hedge fund, bank} \iqlevel{junior}]\label{iq:iv2:markets-and-microstructure:4}
You buy at 100.03 when the mid-price is 100.00. Five minutes later the mid-price is 100.02. Define and compute the effective spread, the realised spread and the price impact, in cents. Which part did the market maker keep?
\end{question}

\begin{question}[$\star$ \iqroles{QT, Trader} \iqfirm{market maker, hedge fund} \iqlevel{junior}]\label{iq:iv2:markets-and-microstructure:5}
A stock has a 1-cent spread. The venue charges a taker 0.3 cents per share and pays a maker 0.2 cents. If you want to buy, how much do you save per share by being filled passively at the bid instead of crossing at the offer? What does a market maker net per share, on a round trip with both sides filled passively, with no adverse selection?
\end{question}

\begin{question}[$\star$ \iqroles{QT, Trader, QR} \iqfirm{market maker, hedge fund} \iqlevel{junior}]\label{iq:iv2:markets-and-microstructure:6}
With a minimum price step of 1 cent, one stock trades at \$5 and another at \$500. What is the smallest spread each can have, in basis points? Which is more likely to sit at a one-tick spread with a long queue, and why? What matters more for trading it, speed or price?
\end{question}

\begin{question}[$\star$ \iqroles{QR, QT} \iqfirm{market maker, hedge fund} \iqlevel{junior}]\label{iq:iv2:markets-and-microstructure:7}
A stock's successive price changes have an autocovariance of $-0.0004$ dollars squared at one lag. Assume the efficient price is a random walk and that each trade is equally likely to be at the bid or the ask, independently. What spread does that imply? Why is the autocovariance negative at all?
\end{question}

\begin{question}[$\star$ \iqroles{QT, Trader, Analyst} \iqfirm{hedge fund, bank} \iqlevel{junior}]\label{iq:iv2:markets-and-microstructure:8}
A day's trades in a stock are 100 shares at 10.00, 200 at 10.20 and 700 at 10.10. What is the volume-weighted average price? You bought 500 shares in the day at an average of 10.08. How did you do against it, in basis points? What does the simple average of the prices give?
\end{question}

\begin{question}[$\star$ \iqroles{QT, Trader, Analyst} \iqfirm{hedge fund, bank} \iqlevel{junior}]\label{iq:iv2:markets-and-microstructure:9}
A stock has 30 million shares sold short, and 2 million shares trade each day. What are its days to cover? If shorts together can be only 20\% of volume when they buy back, how long does it take? Why does this number matter when a stock jumps?
\end{question}

\begin{question}[$\star$ \iqroles{QT, Trader} \iqfirm{market maker, hedge fund, bank} \iqlevel{junior}]\label{iq:iv2:markets-and-microstructure:10}
One American depositary receipt, a US-traded claim, stands for 2 ordinary shares in London. The shares trade at \textsterling 10.00, and one pound buys \$1.25. What is the receipt's fair value? It trades at \$25.20, and all costs of the trade total 30 basis points. What do you do, and how much do you make?
\end{question}

\begin{question}[$\star$ \iqroles{QT, Dev, QD, Trader} \iqfirm{market maker, hedge fund} \iqlevel{junior}]\label{iq:iv2:markets-and-microstructure:11}
Three US venues show offers of 100 shares at 100.01 on $A$, 500 at 100.02 on $B$ and 400 at 100.03 on $C$. You want 300 shares and you like venue $B$. Can you simply send a buy to $B$ at 100.02? What is a trade-through, and what is an intermarket sweep order?
\end{question}

\begin{question}[$\star$ \iqroles{QT, Trader, Dev} \iqfirm{market maker, hedge fund} \iqlevel{junior}]\label{iq:iv2:markets-and-microstructure:12}
An opening auction has buy limit orders of 100 shares at 101, 200 at 100 and 300 at 99, and sell limit orders of 150 at 99, 150 at 100 and 200 at 101. What price maximises the traded volume, and how much trades? Who does not trade?
\end{question}

\begin{question}[$\star$ \iqroles{QT, Trader} \iqfirm{hedge fund, bank} \iqlevel{junior}]\label{iq:iv2:markets-and-microstructure:13}
You must buy 200\,000 shares at 10\% of total volume, your own trades included. Other traders will trade 2 million shares today. How much can you do today, and what share of the others' volume is that? What makes a percentage-of-volume order chase prices?
\end{question}

\begin{question}[$\star$ \iqroles{QT, Trader, Analyst, Risk} \iqfirm{hedge fund, bank} \iqlevel{junior}]\label{iq:iv2:markets-and-microstructure:14}
You want to sell short 10\,000 shares of a US stock today. What must your broker have done before accepting the order? What is a ``fail to deliver'', and is a short sale without a borrow always illegal? Name one kind of participant that has an exception.
\end{question}

\begin{question}[$\star$ \iqroles{QT, QR, Trader} \iqfirm{hedge fund, market maker} \iqlevel{junior}]\label{iq:iv2:markets-and-microstructure:15}
Take this rule: the average cost of buying $Q$ shares, with daily volume $V$ and daily volatility $\sigma$, is about $\sigma\sqrt{Q/V}$. A stock has a daily volatility of 2\%. What does it cost to buy 1\% of a day's volume, and 4\%, in basis points? If you quadruple your order, what happens to the cost in dollars?
\end{question}

\section{Costs, queues and benchmarks}

\begin{question}[$\star\star$ \iqroles{QT, Trader, Analyst} \iqfirm{hedge fund, bank} \iqlevel{junior}]\label{iq:iv2:markets-and-microstructure:16}
You decide to buy 10\,000 shares when the price is \$50. You buy only 8\,000, at an average of \$50.10, and the commission is 1 cent a share. The price closes at \$50.30, and you never buy the other 2\,000. What is the implementation shortfall, in dollars, and as basis points of the paper portfolio? Split it into its parts.
\end{question}

\begin{question}[$\star\star$ \iqroles{QT, QR, Dev} \iqfirm{market maker, hedge fund} \iqlevel{junior, experienced}]\label{iq:iv2:markets-and-microstructure:17}
Your order joins a queue with 20 single-lot orders ahead of it. Events that remove one order from the front of the queue, trades or cancels, arrive as a Poisson process at a rate of 1.5 a second. Treat the removals as coming only from the front. You are filled by the removal that takes your own order, the one after the 20 ahead have gone. How long, on average, until you are filled? What is the standard deviation? What is the chance you are filled within 20 seconds?
\end{question}

\begin{question}[$\star\star$ \iqroles{QT, QR} \iqfirm{market maker, hedge fund} \iqlevel{junior, experienced}]\label{iq:iv2:markets-and-microstructure:18}
A stock has a 1-cent spread, and you post at the bid. If you are first in the queue, small trades fill you, and the price then falls by one tick 20\% of the time. If you are last in the queue, you fill only when something large comes in, and the price then falls by one tick 70\% of the time. If it does not fall, it stays where it is. What do you expect to make per share against the mid-price just before your fill, in each case? How much is the queue position worth?
\end{question}

\begin{question}[$\star\star$ \iqroles{QT, QR, Trader} \iqfirm{market maker, hedge fund} \iqlevel{junior, experienced}]\label{iq:iv2:markets-and-microstructure:19}
A stock is worth either 100 or 120, with a 30\% chance of 120. A quarter of traders know the true value and trade on it. The rest buy or sell with equal probability. Each trades one share, and competition makes you break even on every quote. What are your bid and ask? Why are they not symmetric around the mean? If you see a buy and then a sell, what is your probability of 120?
\end{question}

\begin{question}[$\star\star$ \iqroles{QT, QR} \iqfirm{market maker, hedge fund} \iqlevel{junior, experienced}]\label{iq:iv2:markets-and-microstructure:20}
You make markets and your average buy fills at 10.00 when the mid-price is 10.01. You are paid a 0.2-cent rebate. On average the mid-price is 10.005 one second later and 9.995 one minute later. What are your profits per share marked at those two times? If you do 1 million shares a day, what do you make at each horizon? Which number do you believe, and what should you do?
\end{question}

\begin{question}[$\star\star$ \iqroles{QT, Trader, Analyst} \iqfirm{hedge fund} \iqlevel{junior}]\label{iq:iv2:markets-and-microstructure:21}
You short 100 shares at \$50 for three months. The lender charges a 6\% annual fee on the value, and you earn 4\% a year on the sale proceeds. The stock pays a \$0.50 dividend during that time, and the price is unchanged. What is your cost? How far must the stock fall for you to break even? What is your largest possible loss?
\end{question}

\begin{question}[$\star\star$ \iqroles{QT, Trader, Dev} \iqfirm{hedge fund, market maker} \iqlevel{junior}]\label{iq:iv2:markets-and-microstructure:22}
A US stock closed at \$50 yesterday. Today it trades down to \$44.90 and is now quoted \$44.88 bid and \$44.90 offered. What does the short-sale price test under Rule 201 do? For how long does it hold? At what prices can you now sell short, and can you still hit the bid if you own the shares?
\end{question}

\begin{question}[$\star\star$ \iqroles{QT, Trader, Analyst} \iqfirm{hedge fund, bank} \iqlevel{junior}]\label{iq:iv2:markets-and-microstructure:23}
A trading day has four equal periods. Their shares of the day's volume are 40\%, 15\%, 15\% and 30\%, and the average prices in them are 100.0, 100.5, 101.0 and 100.2. What is the day's volume-weighted average price? You buy on a time-weighted schedule, with equal amounts in each period. What average price do you get, and how far are you from the benchmark in basis points? What would a schedule that follows volume have got?
\end{question}

\begin{question}[$\star\star$ \iqroles{QT, Trader, Analyst} \iqfirm{hedge fund, bank} \iqlevel{junior, experienced}]\label{iq:iv2:markets-and-microstructure:24}
Your order makes up 30\% of a stock's volume today. Others traded at an average of 100.00, and you bought at an average of 100.30, which is what your buying cost you. What is the day's volume-weighted average price? What does your report say against it? What is the true cost, and what happens to the report if you are all of the volume?
\end{question}

\begin{question}[$\star\star$ \iqroles{QT, Trader} \iqfirm{hedge fund, bank} \iqlevel{junior, experienced}]\label{iq:iv2:markets-and-microstructure:25}
You must sell a large block. A dealer will buy it now at 0.8\% below the mid-price. If you work it over two days yourself, the expected cost is 0.5\%, with a standard deviation of 1.5\%. Suppose you penalise a choice by $\lambda$ times its variance, in units of return. At what $\lambda$ are you indifferent? What else would you ask before deciding?
\end{question}

\begin{question}[$\star\star$ \iqroles{QT, QR, LLE, Dev} \iqfirm{market maker, hedge fund} \iqlevel{junior, experienced}]\label{iq:iv2:markets-and-microstructure:26}
A stale offer of 100 shares sits on an exchange, 2 cents below fair value. A market maker's cancel takes a time uniform between 150 and 350 microseconds to arrive, and your order takes a time uniform between 200 and 400, independently. You pay a 0.3-cent taker fee if you trade, and you pay nothing if you miss. What is the chance you win, and what do you expect to make per share per attempt?
\end{question}

\begin{question}[$\star\star$ \iqroles{QT, Trader} \iqfirm{hedge fund, bank} \iqlevel{junior, experienced}]\label{iq:iv2:markets-and-microstructure:27}
You want to buy, and crossing the spread costs 1 cent from the mid-price. You can try a dark pool first, which trades at the mid-price. If it fills you, you save the 1 cent. If it does not, your attempt has leaked and the price rises 0.5 cents before you cross. At what fill probability does it pay to try first? What is the value at a 40\% chance? What is missing from this calculation?
\end{question}

\begin{question}[$\star\star$ \iqroles{QT, Analyst} \iqfirm{market maker, bank} \iqlevel{junior}]\label{iq:iv2:markets-and-microstructure:28}
A dealer fills retail orders in a stock with a 2-cent spread. It gives each order 0.2 cents of improvement on the quote, pays the broker 0.1 cents a share, and on average the price moves 0.2 cents against it afterwards. What does it make per share? Why can it pay for this flow and still beat the exchange's price?
\end{question}

\begin{question}[$\star\star$ \iqroles{QR, QT, DS} \iqfirm{market maker, hedge fund} \iqlevel{junior, experienced}]\label{iq:iv2:markets-and-microstructure:29}
A stock has a true daily volatility of 2\% and a spread of 4 basis points. A day has 390 minutes, and the trade price bounces between bid and ask, independently from one sample to the next. What daily volatility do you estimate if you add up squared one-minute returns? And squared one-second returns? What would you do about it?
\end{question}

\begin{question}[$\star\star$ \iqroles{QT, Trader, Risk} \iqfirm{hedge fund, bank} \iqlevel{junior}]\label{iq:iv2:markets-and-microstructure:30}
You buy \$10 million of a stock with a 30\% annual volatility. Its correlation with an index of 15\% volatility is 0.75. How much of the index do you sell to hedge? What is the remaining annual volatility, in percent and in dollars, and what is that per day, with 16 as the square root of the trading days in a year?
\end{question}

\begin{question}[$\star\star$ \iqroles{QT, QR} \iqfirm{market maker, hedge fund} \iqlevel{junior, experienced}]\label{iq:iv2:markets-and-microstructure:31}
If you quote $\delta$ cents from the mid-price, you are filled at a rate of $e^{-k\delta}$ per second, with $k = 0.5$ per cent. Ignore adverse selection and inventory. What $\delta$ maximises your expected profit rate on one side? What is the fill rate and the profit rate there? How do they change if competition doubles $k$?
\end{question}

\begin{question}[$\star\star$ \iqroles{QT, Trader, QR} \iqfirm{hedge fund, market maker} \iqlevel{junior, experienced}]\label{iq:iv2:markets-and-microstructure:32}
You must buy now or within a minute. The spread is 2 cents, so crossing costs 1 cent from today's mid-price. If you post at the bid instead, you are filled within the minute with probability $p$. If you are not, the mid-price will on average have risen 0.8 cents, and you then cross. Measured against today's mid-price, what do you expect to pay? For what $p$ does it pay to post? Should you also take something off for the mid-price falling when you are filled?
\end{question}

\begin{question}[$\star\star$ \iqroles{QR, QT, DS} \iqfirm{market maker, hedge fund} \iqlevel{junior, experienced}]\label{iq:iv2:markets-and-microstructure:33}
You split a stock's volume into buckets of 100 shares each and classify each bucket's buy and sell volume. Four buckets have buys and sells of 60 and 40, 70 and 30, 50 and 50, and 90 and 10. What is VPIN, the average absolute order imbalance per bucket as a share of bucket size? What is it meant to show? What would you be careful about?
\end{question}

\begin{question}[$\star\star$ \iqroles{QT, Trader, Dev} \iqfirm{hedge fund, market maker} \iqlevel{junior, experienced}]\label{iq:iv2:markets-and-microstructure:34}
A book has 100 shares at each level, 1, 2, 3 and so on cents above the mid-price. You need 1\,000 shares. What does a single sweep cost above the mid-price, in dollars? Alternatively, you buy 100 shares a minute for 10 minutes at one cent above the mid-price, while each slice raises the mid-price 0.2 cents for good. What does that cost? If the stock moves 5 cents a minute in standard deviation, how much risk do the slices take?
\end{question}

\section{Models under questioning}

\begin{question}[$\star\star\star$ \iqroles{QR, QT} \iqfirm{hedge fund, market maker} \iqlevel{junior, experienced}]\label{iq:iv2:markets-and-microstructure:35}
You sell 1\,000 shares in two slices, $x_1$ and then $x_2$. Each slice's own temporary impact costs $a x_i^2$. Your first slice lowers the price the second gets, which costs $b x_1 x_2$. Holding $x_2$ across the period adds a risk penalty $c x_2^2$. Take $a = 0.0002$, $b = 0.0001$ and $c = 0.0001$. What split is best, and what is the total cost? Compare it with an equal split and with doing all of it at once. What happens at $c = 0$, and at $c = 0$ and $b = 2a$?
\end{question}

\begin{question}[$\star\star\star$ \iqroles{QR, QT, Trader} \iqfirm{market maker, hedge fund} \iqlevel{junior, experienced}]\label{iq:iv2:markets-and-microstructure:36}
The best offer has 3 lots, and the best bid has 7. Single-lot buy and sell market orders arrive with equal probability, and there are no new limit orders. The mid-price moves up half a tick if the offer goes first and down half a tick if the bid goes first. What is the chance it moves up? What is the expected move? Compare this with the microprice, and say why the two differ.
\end{question}

\begin{question}[$\star\star\star$ \iqroles{QR, QT} \iqfirm{market maker, hedge fund} \iqlevel{junior, experienced}]\label{iq:iv2:markets-and-microstructure:37}
Trade signs, $+1$ for a buy and $-1$ for a sell, follow a Markov chain, so that the expected next sign is 0.8 times the current one. The price moves only on surprises, by 1 cent for each unit of surprise. How far does the price move on a buy that follows a buy? On a sell that follows a buy? What is the price's total move after a run of five buys, starting from a neutral position? Order flow is strongly autocorrelated: why are price changes not?
\end{question}

\begin{question}[$\star\star\star$ \iqroles{QT, Trader} \iqfirm{market maker, bank} \iqlevel{junior, experienced}]\label{iq:iv2:markets-and-microstructure:38}
A client asks you to buy 1 million shares at once from them. Daily volume is 2 million, and daily volatility is 2\%. You will sell them over time at 20\% of volume. How long does it take? Take your average cost as two thirds of $\sigma\sqrt{Q/V}$. What is it? What is the standard deviation of your outcome from holding a shrinking position? If you charge half of one standard deviation for risk, what discount do you quote? What does hedging with an index do, if the stock's own volatility is 1.2\% a day?
\end{question}

\begin{question}[$\star\star\star$ \iqroles{QR, QT} \iqfirm{market maker, hedge fund} \iqlevel{experienced}]\label{iq:iv2:markets-and-microstructure:39}
You hold $q$ shares of a stock at mid-price $s$, with exponential utility of risk aversion $\gamma$ and a price variance of $\sigma^2$ per unit of time. You will hold until a time $\tau$ from now. Find the price at which you are indifferent to buying one more share, and the same for selling one. What is their average, the reservation price? With $s = 100$, $q = 5$, $\gamma = 0.1$, $\sigma = 2$, $\tau = 0.5$ and a total spread given by $\gamma\sigma^2\tau + (2/\gamma)\ln(1 + \gamma/k)$ with $k = 1.5$, what do you quote?
\end{question}

\begin{question}[$\star\star\star$ \iqroles{QT, Trader, Dev, QD} \iqfirm{market maker, hedge fund} \iqlevel{junior, experienced}]\label{iq:iv2:markets-and-microstructure:40}
A colleague's algorithm puts large bids a few levels below the market, to make others think there is demand, and sells at the offer once the price rises. It cancels the bids before they can fill. A second algorithm cancels 98\% of its orders because its prices keep changing. Are these the same thing? What in US law makes the first one illegal? How would you detect it in your own firm's orders?
\end{question}

\begin{question}[$\star\star\star$ \iqroles{QR, QT, Analyst} \iqfirm{market maker, hedge fund} \iqlevel{experienced}]\label{iq:iv2:markets-and-microstructure:41}
A regulator proposes to halve the minimum price step for a \$20 stock that always trades at a one-tick spread with a long queue. Who gains and who loses: long-term investors, market makers near the front of the queue, market makers near the back, and fast traders who jump ahead? What would you measure afterwards to judge it?
\end{question}

\begin{question}[$\star\star\star$ \iqroles{QR, QT, LLE} \iqfirm{market maker, hedge fund} \iqlevel{junior, experienced}]\label{iq:iv2:markets-and-microstructure:42}
You quote a half-spread of $h$ cents. Ordinary traders hit you at a rate of 1 a second. A public signal moves fair value by 2 cents at a rate of 0.5 a second. There are 3 equally fast firms, one quoting and two who race to pick off the quote; each of the three is equally likely to reach the exchange first. What $h$ makes the quoting firm break even? What does $h$ become with only one firm, and with very many? Would it help to have the race removed?
\end{question}

\begin{question}[$\star\star\star$ \iqroles{QT, QR} \iqfirm{market maker} \iqlevel{junior, experienced}]\label{iq:iv2:markets-and-microstructure:43}
A futures market fills by pure pro-rata. Others show 1\,000 lots at your price. A trade is 200 lots with probability 90\%, after which the price does not move, or 800 lots with probability 10\%, after which it moves 1 tick against you. You earn half a tick against fair value on a fill. If you show $q$ lots, what do you expect to make per trade? What happens as $q$ grows? If you may never take more than 200 lots in one trade, what do you show?
\end{question}

\begin{question}[$\star\star\star$ \iqroles{QR, DS, QT} \iqfirm{hedge fund, bank} \iqlevel{junior, experienced}]\label{iq:iv2:markets-and-microstructure:44}
A fund splits its orders at random between two brokers, 400 to each. The average cost is 3.0 basis points with one and 4.2 with the other, and the standard deviation per order is 30 basis points for both. Is the difference significant? How many orders per broker do you need to detect a difference of 1.2 basis points with 80\% power at the 5\% level, two-sided? What does the cost's size have to do with how you should run this?
\end{question}

\begin{question}[$\star\star\star$ \iqroles{QR, QT} \iqfirm{hedge fund, market maker} \iqlevel{junior, experienced}]\label{iq:iv2:markets-and-microstructure:45}
You sell 50\,000 shares in each of two periods. Each share sold lowers the price permanently by $b = \$0.000001$, and trades happen at the price after that period's selling. A predator who knows your plan sells $y$ shares in the first period alongside you and buys them back after your second period, when that buying raises the price back by $by$. What $y$ maximises the predator's profit, and what is the profit? What is your extra cost, and as a share of your cost with no predator? What does this say about how you should schedule?
\end{question}

\begin{question}[$\star\star\star$ \iqroles{QR, QT, Trader} \iqfirm{hedge fund, bank} \iqlevel{experienced}]\label{iq:iv2:markets-and-microstructure:46}
You must sell a quarter of a day's volume, with a daily volatility of 2\%. Suppose that over a horizon of $T$ days your expected cost is $\tfrac23\sigma\sqrt{Q/(VT)}$ and that your risk is $\sigma\sqrt{T/3}$, and you minimise the cost plus one unit of the risk. What horizon is best? What are the two terms there? What does the empirical square-root law say that makes you wary of this cost model?
\end{question}

\begin{question}[$\star\star\star$ \iqroles{QT, Trader, Risk} \iqfirm{market maker, hedge fund} \iqlevel{junior, experienced}]\label{iq:iv2:markets-and-microstructure:47}
A US stock's reference price is \$100, in the middle of the trading day, and its limit up-limit down band is 5\%. At what prices do the bands sit? It trades down to the lower band, and offers stay there with bids below it. What happens after 15 seconds? How does the band change in the first and last parts of the day? You have sold a put on it and are delta-hedged: what is your problem during the pause, and what do you do?
\end{question}

\begin{question}[$\star\star\star\star$ \iqroles{QR} \iqfirm{hedge fund, market maker} \iqlevel{experienced}]\label{iq:iv2:markets-and-microstructure:48}
An asset's value $v$ is normal around $p_0$ with variance $\sigma_v^2$. A single insider sees $v$ and submits an order $x$. Noise traders submit $u$, normal with mean zero and variance $\sigma_u^2$, independent of $v$. A market maker sees only $x + u$ and sets the price at $p_0 + \lambda(x + u)$, equal to the expected value given the flow. Find the insider's strategy $x = \beta(v - p_0)$ and $\lambda$. With $\sigma_v = 2$ dollars and $\sigma_u = 1\,000$ shares, what are $\beta$, $\lambda$, the insider's expected profit, and the price's remaining variance? Who pays the insider?
\end{question}

\begin{question}[$\star\star\star\star$ \iqroles{QR} \iqfirm{hedge fund, bank, market maker} \iqlevel{experienced}]\label{iq:iv2:markets-and-microstructure:49}
You must sell $X$ shares by time $T$, holding $x(t)$ shares at time $t$. Selling at a rate $v$ costs you $\eta v$ per share. You minimise $\int_0^T \eta v^2\,dt + \lambda\sigma^2\int_0^T x^2\,dt$. Derive the optimal path. Let $\kappa^2 = \lambda\sigma^2/\eta$. What share is sold by half way if $\kappa T = 1$, and if $\kappa T = 3$? In each case, how do the cost and the price variance compare with selling at a constant rate? What is the path when $\lambda = 0$?
\end{question}

\begin{question}[$\star\star\star\star$ \iqroles{QR} \iqfirm{hedge fund, market maker} \iqlevel{experienced}]\label{iq:iv2:markets-and-microstructure:50}
You sell $X$ shares in three trades at equally spaced times, of sizes $x_0$, $x_1$ and $x_2$. A trade of size $x$ moves the price by $\kappa x$, costing you $\tfrac12\kappa x^2$ on itself. The earlier move then decays, and a fraction $a$ of it is still there at the next trade time. So the total cost is $\tfrac12\kappa\left(x_0^2 + x_1^2 + x_2^2\right) + \kappa a\left(x_0x_1 + x_1x_2\right) + \kappa a^2 x_0x_2$. Find the best sizes. With $a = 0.5$, what are the shares, and what do you save against equal thirds, and against all at once? What happens as $a$ goes to 0 and to 1?
\end{question}
